\documentclass{beamer}
\usepackage[utf8]{inputenc}
\usepackage[brazilian]{babel}
\usepackage{booktabs}
\usepackage{xcolor}

% Tema CERISE — mesma identidade visual dos outros slides do projeto.
\definecolor{cerisemagenta}{HTML}{A6186B}
\definecolor{ceriserosa}{HTML}{E6186D}
\definecolor{cerisecoral}{HTML}{F5642D}
\definecolor{ceriseroxo}{HTML}{5B1A6B}

\setbeamercolor{title}{fg=cerisemagenta}
\setbeamercolor{frametitle}{fg=white,bg=cerisemagenta}
\setbeamercolor{structure}{fg=cerisemagenta}
\setbeamercolor{block title}{fg=white,bg=ceriseroxo}
\setbeamercolor{block body}{bg=cerisemagenta!5}
\setbeamercolor{alerted text}{fg=cerisecoral}
\setbeamercolor{item}{fg=cerisemagenta}

\newenvironment{numerobloco}[1]{\begin{block}{#1}}{\end{block}}
\newenvironment{decisaobloco}[1]{\begin{alertblock}{#1}}{\end{alertblock}}

\title{Progresso do Harbor}
\subtitle{Trace, agentes, RAG multimodal e migração LangGraph}
\author{Yan Santos Leite}
\date{8--9 de setembro de 2026}

\begin{document}

\frame{\titlepage}

\begin{frame}{Resumo}
  \begin{itemize}
    \item Quatro frentes fechadas neste período: \textbf{instrumentação}
      (trace + custo/versão), \textbf{agentes ReAct} (LangGraph, com trace, sobre
      duas tarefas de benchmark), \textbf{RAG multimodal} (baseline medido, benchmark
      de retrieval visual), e a \textbf{migração para LangGraph} do self-repair e do
      roteador (produção).
    \item Dois resultados negativos medidos e \textbf{não promovidos} — Agentic RAG e
      Contextual Retrieval — tratados aqui como resultado legítimo, não falha.
    \item Três bugs reais de import encontrados só ao testar o dashboard ao vivo,
      corrigidos e confirmados.
  \end{itemize}
\end{frame}

\begin{frame}{Instrumentação: trace e custo (Fases 1--3 da disciplina PDC)}
  \begin{numerobloco}{O que existe agora}
    \begin{itemize}
      \item \texttt{shared/trace.py} — esquema de trace espelhando \texttt{gen\_ai.*}
        (OpenTelemetry) + camada semântica de proveniência (evidências, relações
        SUPPORT/CONTRADICT).
      \item \texttt{shared/ollama\_client.py} — captura tokens de entrada/saída e
        modelo em toda chamada Ollama (antes: 13 pontos descartavam essa informação).
      \item \texttt{shared/rollout\_card.py} — relatório de replicação em Markdown a
        partir dos traces.
    \end{itemize}
  \end{numerobloco}

  \pause
  \textbf{Por quê:} sem isso, nenhuma comparação de custo (LangGraph vs. Python puro,
  qwen2.5:7b vs. DeepSeek) seria possível com número real — só suposição.
\end{frame}

\begin{frame}{Agente ReAct em LangGraph, sobre duas tarefas}
  \small
  \begin{numerobloco}{Resultado real (Rollout Cards)}
    \begin{center}
    \begin{tabular}{lccc}
      \toprule
      \textbf{Tarefa} & \textbf{Execuções} & \textbf{Sucesso} & \textbf{Tokens (in/out)} \\
      \midrule
      SQL (self-repair)     & 3 & 1/3 (33\%) & 27.163 / 1.342 \\
      HotpotQA (replicação) & 2 & 1/2 (50\%) & — \\
      \bottomrule
    \end{tabular}
    \end{center}
    \footnotesize Fonte: Rollout Cards em \texttt{eval/traces/}. Amostra pequena —
    valida o padrão de trace, não é conclusão estatística sobre o modelo.
  \end{numerobloco}

  \pause
  \textbf{O que isso significou:} 2 das 3 execuções SQL travaram no MESMO erro de
  sintaxe repetido, sem o modelo usar a Observation de erro para se corrigir.

  \pause
  \begin{decisaobloco}{O que decidimos}
    Entregar sempre os três artefatos (código, trace, Rollout Card) — regra permanente
    na skill \texttt{padrao-react-raciocinio-acao}.
  \end{decisaobloco}
\end{frame}

\begin{frame}{RAG multimodal: baseline e benchmark de retrieval visual}
  \small
  \begin{numerobloco}{O que medimos}
    Caption-then-embed (baseline): Recall@3 86\%, MRR 0,524 (corpus próprio, 4 imagens).
    Benchmark 2 estágios (retrieval + rerank via ColModernVBERT) sobre ViDoRe (subset
    real) e corpus próprio, comparando top-k vs. top-p.
  \end{numerobloco}

  \pause
  \textbf{O que isso significou:} no corpus maior (ViDoRe), top-p empatou com "sem
  rerank" a 1/3 do custo. No corpus pequeno (Harbor), top-p piorou muito — investigado
  a fundo: não era limiar mal calibrado, era o 1º estágio nunca gerar score real
  (\texttt{usar\_rerank=False} sempre devolve \texttt{score=0.0}).

  \pause
  \begin{decisaobloco}{O que decidimos}
    Não decidir arquitetura de produção a partir de um número que não mede o que deveria
    medir — item de continuidade: dar ao 1º estágio um sinal de score real antes de
    re-calibrar.
  \end{decisaobloco}
\end{frame}

\begin{frame}{Migração LangGraph: self-repair e roteador promovidos}
  \small
  \begin{numerobloco}{O que medimos (harness completo, 67 perguntas)}
    \begin{center}
    \begin{tabular}{lcc}
      \toprule
      \textbf{Métrica} & \textbf{Antes} & \textbf{Depois} \\
      \midrule
      Roteamento     & 56/67 & 55/67$^*$ \\
      Faithfulness   & 64\%  & 64\%     \\
      Divergências de rota (roteador) & — & 0/67 \\
      \bottomrule
    \end{tabular}
    \end{center}
    \footnotesize $^*$ instabilidade pré-existente do desempate por LLM numa pergunta
    de fronteira — confirmada também no roteador original, não regressão.
  \end{numerobloco}

  \pause
  \begin{decisaobloco}{O que decidimos}
    Ambos promovidos para produção. Confirmado funcionando no dashboard real, incluindo
    a correção de 3 bugs de import descobertos só nesse teste manual.
  \end{decisaobloco}
\end{frame}

\begin{frame}{Dois resultados negativos, medidos e não promovidos}
  \begin{numerobloco}{Agentic RAG}
    Recall@1: 12/15 (chamada única) $\to$ \alert{11/15} (com critic node). Causa: o
    chunk certo às vezes não está no topo do documento certo — reformular a query não
    resolve um problema de chunking.
  \end{numerobloco}
  \begin{numerobloco}{Contextual Retrieval}
    Precision@5: 50,0\% $\to$ \alert{46,4\%}, MRR levemente melhor (+0,008). Manuais
    curtos (40--181 linhas) não se beneficiam como a literatura original (documentos
    longos) prevê.
  \end{numerobloco}

  \pause
  \textbf{Por que isso está aqui:} medir antes de decidir funcionou como deveria —
  evitou promover duas mudanças que pioram o sistema.
\end{frame}

\begin{frame}{Próximos passos}
  \begin{itemize}
    \item Diagnóstico em camadas e NL-to-SQL (geração) — próximos itens do roadmap
      \texttt{migrar-para-langchain}, maior risco.
    \item Dar ao 1º estágio do benchmark multimodal um score real antes de re-calibrar
      top-p.
    \item Instrumentar trace no dashboard/chat de produção real (hoje só nos agentes
      de teste).
    \item Checkpointer real em \texttt{agents/react.py}, harness formal de avaliação de
      agente, multiagente Supervisor + especialistas (do slide CERISE).
  \end{itemize}
\end{frame}

\end{document}
